\chapter{Методика экспериментов}

Каждый механизм контура сравнивается с вариантом без него на одном и том же
отказе, при постоянной нагрузке и с тремя повторами на конфигурацию. Прогон
полностью автоматизирован. Целевые показатели и гипотезы зафиксированы до
проведения экспериментов и в ходе работы не менялись.

\section{Целевые показатели}

Работа считается выполненной при достижении значений, приведённых в
табл.~\ref{tab:targets}.

\begin{table}[ht]
  \centering
  \caption{Целевые показатели}
  \label{tab:targets}
  \small
  \begin{tabularx}{\textwidth}{@{}lXl@{}}
    \toprule
    Область & Показатель & Целевое значение \\
    \midrule
    SLO & Доля успешных запросов в штатном режиме & $\ge 99.0\%$ \\
    SLO & Задержка p95 при нагрузке 50 запросов/с & $< 200$ мс \\
    SLO & Скорость расходования бюджета ошибок & отслеживается, алерт \\
    Восстановление & Время обнаружения отказа (TTD) & $< 60$ с \\
    Восстановление & Время восстановления (RTO) & $< 5$ мин \\
    Восстановление & Потерянные запросы за инцидент & фиксируются \\
    Политики & Размещённые нарушители после включения контроля & 0 \\
    Планировщик & Снижение доли ошибок относительно наивного выбора
      в сценарии деградации & $\ge 20\%$ \\
    \bottomrule
  \end{tabularx}
\end{table}

Показатель допустимой потери данных (RPO) для сервиса без состояния не
определён; вместо него фиксируется число запросов, потерянных за время
инцидента.

\section{Гипотезы и правило принятия решения}

Проверяются пять гипотез:
\begin{enumerate}[label=H\arabic*.]
  \item В штатном режиме доступность не ниже 99.0\% и задержка p95 менее
        200~мс; алерт по скорости расходования бюджета ошибок срабатывает
        при отказе.
  \item Контроллер аварийного восстановления обеспечивает TTD менее 60~с,
        RTO менее 5~мин и меньшее число потерянных запросов, чем режим без
        переключения трафика.
  \item При включённом контроле политик не размещается ни один нарушитель;
        в режиме аудита нарушители размещаются и фиксируются в каталоге.
  \item Доля ошибок потока нагрузок при разрабатываемом планировщике ниже,
        чем при наивном, не менее чем на 20\%.
  \item Вывод гипотезы H4 устойчив к выбору весов: при любых весах с
        $w_1 + w_2 > 0$ доля размещений в деградировавший кластер ниже, чем
        у наивного планировщика.
\end{enumerate}

Гипотеза считается подтверждённой, если показатель достигнут на среднем
трёх повторов и эффект имеет тот же знак в каждом повторе, а для парных
сравнений~--- в каждой паре. Пороговые цели (отсутствие нарушителей, TTD
менее 60~с) проверяются в каждом прогоне. Составная гипотеза подтверждена,
только если подтверждены все её части; в противном случае она считается
подтверждённой частично. Статистические тесты не применяются: при трёх
повторах они неинформативны, и вместо них используется требование
одинакового знака эффекта в каждой паре.

\section{Нагрузка и структура прогона}

Нагрузку создаёт k6~\cite{k6} с постоянной частотой 50 запросов в секунду.
Постоянная частота выбрана вместо постепенного наращивания, чтобы нагрузка не
стала вторым источником изменения метрик. Каждый запрос записывается с
отметкой времени; по этим записям вычисляются потерянные запросы и доля
ошибок на стороне клиента.

Прогон состоит из 3~минут базовой линии, 5~минут отказа и 5~минут
восстановления. На каждую конфигурацию выполняется три повтора, причём
сравниваемые режимы чередуются, чтобы постепенный дрейф стенда не
приходился на один режим. Параметры, сырые ряды метрик, записи каталога и
сводка каждого прогона сохраняются в репозитории.

\section{Внесение отказов}

Отказы вносятся с помощью Chaos Mesh~\cite{chaos-mesh}. Время внесения и
снятия отказа берётся из событий самого эксперимента Chaos Mesh, а не из
наблюдений контура.

\section{Сценарии}

Используются три сценария (табл.~\ref{tab:scenarios}).

\begin{table}[ht]
  \centering
  \caption{Сценарии экспериментов}
  \label{tab:scenarios}
  \small
  \begin{tabularx}{\textwidth}{@{}p{0.17\textwidth}XXX@{}}
    \toprule
    Сценарий & Отказ & Сравнение & Измеряется \\
    \midrule
    1. Отказ кластера & Отказ всех подов сервиса в кластере B на 5~мин;
      нагрузка на долгоживущий сервис & Контроллер восстановления
      включён и в режиме наблюдения & TTD, RTO, потерянные запросы,
      доступность \\
    2. Деградация сети & Задержка 500~мс и потеря 10\% пакетов до backend
      кластера B; нагрузка на поток и фоновые 5 запросов/с на
      долгоживущий сервис & Разрабатываемый и наивный планировщики &
      Доля ошибок потока в окне отказа, реакция планировщика \\
    3. Нарушение политики & 5 попыток разместить в обход планировщика
      нагрузку, требующую отсутствующей у кластера сертификации &
      Политики в режиме контроля и в режиме аудита & Размещённые
      нарушители, записи в журнале, задержка фиксации \\
    \bottomrule
  \end{tabularx}
\end{table}

В сценарии~2 контроллер восстановления работает в режиме наблюдения: если бы
он переключал трафик, эффект планировщика смешался бы с эффектом
восстановления. Фоновая нагрузка на долгоживущий сервис нужна, чтобы у
каждого кластера были метрики: кластер без трафика получает нейтральные
оценки и проигрывает кластеру с трафиком при любом его состоянии.

\section{Метрики и свёртка повторов}

\begin{itemize}
  \item TTD~--- интервал от внесения отказа до перехода кластера в состояние
        деградации; RTO~--- интервал от внесения отказа до завершения
        переключения трафика. Оба вычисляются по записи инцидента в каталоге.
  \item Потерянные запросы~--- неуспешные запросы клиента от внесения отказа
        до переключения трафика, а без переключения~--- за всё время отказа.
  \item Снижение доли ошибок в сценарии~2:
        \[
          \Delta = 1 - \frac{\overline{ER}_{\mathrm{meta}}}{\overline{ER}_{\mathrm{naive}}},
        \]
        где $\overline{ER}$~--- средняя по повторам доля ошибок потока на
        стороне клиента в окне от внесения до снятия отказа. Серверный ряд
        из Thanos используется для проверки.
  \item Реакция планировщика~--- доля размещений в деградировавший кластер
        за окно отказа, время до снижения оценки этого кластера в каталоге и
        время до первого решения в обход него.
\end{itemize}

\section{Критерии отбраковки прогонов}

Прогон признаётся недействительным и повторяется, если нарушены условия
сравнения: отсутствует событие снятия отказа, адреса узлов разошлись с
маршрутами шлюза или перед началом прогона каталог не видел оба кластера
в исправном состоянии. Признаки отбраковки записываются в сводку прогона.
